\documentclass{article}

\usepackage{amsmath,amssymb}
\usepackage{xurl}
\usepackage[hidelinks]{hyperref}
\setlength{\emergencystretch}{2em}

% Shared commands for notes in docs/paper-gaps/.

\DeclareUnicodeCharacter{2080}{\ifmmode{}_{0}\else\textsubscript{0}\fi}
\DeclareUnicodeCharacter{2081}{\ifmmode{}_{1}\else\textsubscript{1}\fi}
\DeclareUnicodeCharacter{2082}{\ifmmode{}_{2}\else\textsubscript{2}\fi}
\DeclareUnicodeCharacter{2099}{\ifmmode{}_{n}\else\textsubscript{n}\fi}

\newcommand{\Lean}{\textsc{Lean}}
\ExplSyntaxOn
\NewDocumentCommand{\leanid}{m}{
  \ifmmode
    \textsf{\detokenize{#1}}
  \else
    \group_begin:
    \urlstyle{sf}
    \tl_set:Nn \l_tmpa_tl {#1}
    \tl_replace_all:Nnn \l_tmpa_tl {\_} {_}
    \exp_args:NV \path \l_tmpa_tl
    \group_end:
  \fi
}
\ExplSyntaxOff
\providecommand{\tr}{\operatorname{tr}}
\newcommand{\tnleanIssueBase}{https://github.com/LionSR/TNLean/issues/}
\newcommand{\tnleanPrBase}{https://github.com/LionSR/TNLean/pull/}
\newcommand{\tnleanDocBase}{https://sirui-lu.com/TNLean/docs/}
\newcommand{\ghissue}[1]{\href{\tnleanIssueBase#1}{GitHub issue \##1}}
\newcommand{\ghpr}[1]{\href{\tnleanPrBase#1}{PR \##1}}
\newcommand{\doclink}[1]{\href{\tnleanDocBase#1.html}{\path{#1}}}

% Dirac notation and the identity operator, as in blueprint/src/macros/common.tex.
% \providecommand keeps note-local definitions authoritative.
\usepackage{dsfont}
\providecommand{\ket}[1]{|#1\rangle}
\providecommand{\bra}[1]{\langle #1|}
\providecommand{\braket}[2]{\langle #1 | #2 \rangle}
\providecommand{\ketbra}[2]{|#1\rangle\!\langle #2|}
\providecommand{\Id}{\mathds{1}}
\providecommand{\one}{\mathds{1}}
\providecommand{\C}{\mathbb{C}}
\providecommand{\MN}[1]{M_{#1}(\C)}
\providecommand{\MD}{M_D(\C)}

% External referencing for paper sources.  The build helper writes sanitized
% auxiliary files under build/paper-aux/ and excludes duplicated source labels.
\usepackage{xr}

% Import a generated paper auxiliary file with a caller-supplied label prefix.
% The candidate paths support builds from docs/paper-gaps/, the repository root,
% and build/paper-aux/ itself.
\newcommand{\importpaperaux}[2]{%
  \IfFileExists{../../build/paper-aux/#2.aux}{%
    \externaldocument[#1]{../../build/paper-aux/#2}%
  }{%
    \IfFileExists{build/paper-aux/#2.aux}{%
      \externaldocument[#1]{build/paper-aux/#2}%
    }{%
      \IfFileExists{#2.aux}{%
        \externaldocument[#1]{#2}%
      }{}%
    }%
  }%
}

% General external reference macros
\newcommand{\paperref}[2]{\ref{#1-#2}}
\newcommand{\papereqref}[2]{(\ref{#1-#2})}

% CPSV16 source (1606.00608 / MPDO-22-12-17-2.tex) external document and shortcuts
\importpaperaux{cpsv16-}{1606.00608/MPDO-22-12-17-2-tnlean}
\newcommand{\cpsvref}[1]{\paperref{cpsv16}{#1}}
\newcommand{\cpsveqref}[1]{\papereqref{cpsv16}{#1}}

% Verdict marker: \gapnote{<kind>}{<status>}. Kind is the policy
% classification (clarification, local-correction, scope-restriction,
% unfaithful, false-source, open-gap); status is open, wip (elimination
% underway), resolved, or historical. Machine-read by the site index;
% prints a verdict line.
\newcommand{\gapnote}[2]{\par\noindent\textbf{Verdict:} #1 (#2).\par}


\title{The Intertwiner Count in the Quadratic Reconstruction Argument}
\author{}
\date{2026-08-22}

\begin{document}
\maketitle
\gapnote{local-correction}{open}

\section{The printed dependence step}

The proof of Theorem~7 of P\'erez-Garc\'ia, Verstraete, Wolf, and Cirac
\cite{PerezGarcia2007Matrix} compares two translation-invariant matrix product
representations through matrices acting on a $D^2$-dimensional doubled bond
space. The local source is
\path{Papers/quant-ph_0608197/MPSarchive.tex}, lines~1080--1095. It produces
$D^4$ invertible matrices $W_1,\ldots,W_{D^4}$ satisfying
\begin{align}
    W_k(C_i\otimes I) & =(B_i\otimes I)W_{k+1}
                      &                        & (1\leq k\leq D^4-1),                                      \label{eq:source-chain}
\end{align}
and then chooses an index $n$ for which $W_1,\ldots,W_{n-1}$ are linearly
independent but $W_n$ belongs to their span.

The ambient matrix space is
\begin{align}
    M_{D^2}(\mathbb C),\qquad \dim_{\mathbb C}M_{D^2}(\mathbb C)=D^4.
\end{align}
Consequently, $D^4$ matrices need not be linearly dependent. The dependence
asserted after~\eqref{eq:source-chain} does not follow from the printed count.

\section{The corrected cut count}

A window of $D^4$ consecutive open-boundary matrices has $D^4+1$ cuts: one
before the first site, one after the last site, and $D^4-1$ internal cuts. If a
comparison of the two open-boundary representations supplies an intertwiner at
every cut, then one obtains
\begin{align}
    W_0,W_1,\ldots,W_{D^4} & \in M_{D^2}(\mathbb C),                 \label{eq:full-chain}                                                                                       \\
    W_k(C_i\otimes I)      & =(B_i\otimes I)W_{k+1}
                           &                                                               & (0\leq k\leq D^4-1).                                      \label{eq:full-relations}
\end{align}
The $D^4+1$ vectors in~\eqref{eq:full-chain} are necessarily linearly
dependent as a family. Since every $W_k$ is invertible, $W_0\neq0$. Choose the
first dependent vector. Then there is an index $1\leq n\leq D^4$ such that
$W_0,\ldots,W_{n-1}$ are linearly independent and $W_n$ lies in their span.
After shifting indices, this is exactly the dependence hypothesis required by
the chain-combination lemma at source lines~1022--1049.

The current formal quadratic-system structure stores only gauges attached to
the $D^4$ sites. It therefore constructs the $D^4-1$ internal-cut
intertwiners and proves the $D^4-2$ relations between consecutive internal
cuts. It has no fields representing the gauges at the two endpoint cuts.
Those endpoint intertwiners must be derived from a future comparison theorem
for the surrounding open-boundary representations. They are not consequences
of the internal quadratic equations alone.

\section{Verdict}

The source's $D^4$-intertwiner count is off by one for the stated dimension
argument. The local correction is to use the full chain of $D^4+1$ cut
intertwiners around the $D^4$-site window. Once the two endpoint cuts and all
$D^4$ adjacent relations are supplied, linear dependence and the subsequent
combination argument are valid. The formal development currently proves only
the internal chain and does not claim that this chain forces dependence, so
the reconstruction capstone remains open.

\bibliographystyle{alpha}
\bibliography{references}

\end{document}
